\section{Proofs of the endpoint equation and gap rate}
\label{v12:app:mellin-proofs}

\subsection{The explicit local models and all localization errors}

In coordinates $f_1(t)=u(t)$, $f_2(t)=u(t-1)$, $0<t<1$, the finite
Hilbert transform has kernel
\[
 \frac1\pi\begin{pmatrix}
 \operatorname{pv}(t-s)^{-1}&(t-s+1)^{-1}\\
 (t-s-1)^{-1}&\operatorname{pv}(t-s)^{-1}
 \end{pmatrix}.
\]
On $L^2(0,\infty)$ let $H_\infty,C_\infty$ have kernels
$[\pi(t-s)]^{-1}$, in principal value, and $[\pi(t+s)]^{-1}$.
On four copies of this space set
\[
 H_e=\begin{pmatrix}
 H_\infty&0&0&C_\infty\\0&H_\infty&0&0\\
 0&0&-H_\infty&0\\-C_\infty&0&0&-H_\infty
 \end{pmatrix},\quad Q_e=\tfrac12(I+iH_e),\quad
 J_e=\operatorname{diag}(J_\theta,J_\theta),\quad
 M_\lambda=Q_e-\lambda J_e.
\]
For the unitary Mellin transform
$\mathcal Mf(\omega)=(2\pi)^{-1/2}\int_0^\infty
t^{-1/2-i\omega}f(t)\,dt$, the two scalar multipliers are
$-i\tanh\pi\omega$ and $\operatorname{sech}\pi\omega$.
Indeed these follow from the principal-value cotangent integral and
the beta integral at exponent $1/2+i\omega$.  Thus $M_\lambda$ has
multiplier \eqref{v12:eq:endpoint-symbol}.  Expansion of its determinant
gives $\lambda^4-\lambda^2(q^2+r^2)+q^2r^2-c^2qr$; using $c^2=qr$
proves the asserted formula.  The $2\times2$ block of $Q_e$ on
channels 1 and 4 is a projection, and its other entries are $q,r$.
Consequently $0\le Q_e\le I$.  The inverse multiplier is bounded on
compact subsets of $\Omega$: its denominator is bounded away from zero,
and its adjugate is bounded.

Choose $0<\delta<1/4$ and a real smooth cutoff $0\le\chi\le1$, equal to one on
$[0,\delta]$ and zero on $[2\delta,\infty)$, such that
$\psi(t)=\sqrt{1-\chi(t)^2-\chi(1-t)^2}$ is smooth.  A smooth
sine/cosine transition has this property.  Extend $\psi$ by zero.
Define
\[
 Ef(t)=\chi(t)(f_1(t),f_2(t),f_1(1-t),f_2(1-t))^T,
 \qquad \mathcal If=\psi f,
\]
where the outputs are in $L^2((0,\infty);\mathbb C^4)$ and
$L^2(\mathbb R;\mathbb C^2)$, respectively, with zero extension.
Their adjoints are
\[
 E^*v(t)=\begin{pmatrix}
 \chi(t)v_1(t)+\chi(1-t)v_3(1-t)\\
 \chi(t)v_2(t)+\chi(1-t)v_4(1-t)
 \end{pmatrix},\qquad \mathcal I^*w(t)=\psi(t)w(t).
\]
Hence $E^*E+\mathcal I^*\mathcal I=1$, $EJ_\theta=J_eE$, and $\mathcal IJ_\theta=J_\theta\mathcal I$.
Let $Q_i$ be the whole-line Cauchy projection on each of the two interior
components and $N_\lambda=Q_i-\lambda J_\theta$.  In angular Fourier
frequency,
\begin{equation}\label{v12:eq:interior-inverse}
 n_\lambda(\omega)^{-1}=
 \begin{cases}-\lambda^{-1}J_\theta,&\omega<0,\\
 (I+\lambda J_\theta)/(1-\lambda^2),&\omega>0.
 \end{cases}
\end{equation}

The errors $\Delta_e=EQ_0-Q_eE$ and $\Delta_i=\mathcal IQ_0-Q_i\mathcal I$ are
independent of $\lambda,\theta$.  For $t>0$, $0<s<1$, the first has
kernel
\begin{equation}\label{v12:eq:endpoint-defect-kernel}
 \frac{i}{2\pi}\begin{pmatrix}
 \dfrac{\chi(t)-\chi(s)}{t-s}&
 \dfrac{\chi(t)-\chi(1-s)}{t+1-s}\\[5pt]
 \dfrac{\chi(t)}{t-1-s}&
 \dfrac{\chi(t)-\chi(s)}{t-s}\\[5pt]
 \dfrac{\chi(t)-\chi(1-s)}{1-t-s}&
 \dfrac{\chi(t)}{2-t-s}\\[5pt]
 -\dfrac{\chi(t)-\chi(s)}{t+s}&
 \dfrac{\chi(t)-\chi(1-s)}{1-t-s}
 \end{pmatrix}.
\end{equation}
For $t\in\mathbb R$, $0<s<1$, the second has kernel
\begin{equation}\label{v12:eq:interior-defect-kernel}
 \frac{i}{2\pi}\begin{pmatrix}
 (\psi(t)-\psi(s))/(t-s)&\psi(t)/(t-s+1)\\
 \psi(t)/(t-s-1)&(\psi(t)-\psi(s))/(t-s)
 \end{pmatrix}.
\end{equation}
Divided differences take their continuous values; a term supported where
its numerator is identically zero is zero.  The kernels are square integrable: diagonal quotients are bounded by
cutoff derivatives; corner quotients vanish below distance $\delta$,
with denominator at least $\delta$ elsewhere; unmatched cross terms
are separated from their poles. Output tails are $O(|t|^{-1})$.
The touching-endpoint Carleman singularity remains in the model.

\subsection{The scalar equation and its multiplicities}

\begin{proof}[Proof of Theorem~\ref{v12:thm:explicit-gap-equation}]
The intertwining errors satisfy
$EP_\lambda=M_\lambda E+\Delta_e$ and
$\mathcal IP_\lambda=N_\lambda\mathcal I+\Delta_i$.  Therefore
\[
 EP_\lambda^2=M_\lambda^2E+M_\lambda\Delta_e+\Delta_eP_\lambda,
\]
with the analogous interior identity.  Multiplying by the inverse
squares and using $E^*E+\mathcal I^*\mathcal I=1$ gives
\eqref{v12:eq:explicit-preconditioner}.  All inverse powers here are
algebraic powers, not adjoint products.  The explicit multipliers are
bounded and holomorphic on $\Omega$, so $K_\lambda$ is holomorphic in
$\mathfrak S_2$.

For real gap $\lambda$, self-adjoint model pencils of norm at most
$1+|\lambda|$ have inverse squares at least $(1+|\lambda|)^{-2}I$.
Thus $R_\lambda$ is invertible and $P_\lambda^2$ is Fredholm, with
$\ker P_\lambda^2=\ker P_\lambda$. The exact frequency correspondence
proves the real zero equivalence.

At a real isolated eigenvalue $\alpha$ of multiplicity $k$, the
bounded-product correspondence identifies the nonzero generalized
eigenspace of $J_\theta Q_0$ with that of the self-adjoint physical
operator.  It is therefore semisimple of dimension $k$.  Its Riesz
decomposition splits $J_\theta Q_0-\lambda$ into
$(\alpha-\lambda)I_k$ and an invertible analytic complement.  The
analytic Fredholm pencil $P_\lambda=J_\theta(J_\theta Q_0-\lambda)$
has characteristic multiplicity $k$; its square has multiplicity $2k$.
Multiplication by the locally invertible analytic $R_\lambda$ preserves
this number.  Finite-dimensional Schur reduction of an identity plus
$\mathfrak S_2$ family shows that its regularized determinant has this
zero order; the regularizing factor is a nonvanishing exponential.
\end{proof}

The equation permits finite-rank error control.  For $T\ge2$, the
displayed kernels give
\[
 \|\mathbf1_{t>T}\Delta_e\|_2^2\le\frac3{2\pi^2(T-1)},\qquad
 \|\mathbf1_{|t|>T}\Delta_i\|_2^2\le\frac1{\pi^2(T-1)}.
\]
There are six nonzero endpoint tail entries, bounded by
$(2\pi)^{-1}(t-1)^{-1}$, and two diagonal interior entries on each tail.
On the remaining finite rectangles, midpoint step approximations with
mesh widths $h_t,h_s$ have Hilbert--Schmidt error at most
$\sqrt T(B_th_t+B_sh_s)/2$, or
$\sqrt{2T}(B_th_t+B_sh_s)/2$ for the interior, where $B_t,B_s$ bound
the Frobenius norms of the two kernel derivatives.  Such bounds follow
from cutoff second derivatives and the canceled corner formulas.
The squared rectangle and tail errors add.

Replacing only the errors in \eqref{v12:eq:explicit-preconditioner}
by these finite-rank approximations gives, for example,
\[
 \|K_e-K_{e,n}\|_2\le
 (\|M_\lambda^{-1}\|^2(1+|\lambda|)+\|M_\lambda^{-1}\|)
                    \|\Delta_e-\Delta_{e,n}\|_2.
\]
Every multiplier is explicit; matrix evaluation errors must also be
controlled in a certificate.  A contour may certify zeros only when
it bounds a domain $D$ with $\overline D\Subset\Omega$ and
$R_\lambda$ proved invertible throughout $\overline D$.  If model
inverse norms are at most $b$ on a disk about real $\lambda_0$, then
$\|R_\lambda-R_{\lambda_0}\|\le2b^3|\lambda-\lambda_0|$; a disk with
$2b^3|\lambda-\lambda_0|<(1+|\lambda_0|)^{-2}$ suffices.
On the contour, a verified bound
$\|(I+K_{\lambda,n})^{-1}\|\,\|K_\lambda-K_{\lambda,n}\|<1$
keeps the interpolation invertible.  Half the finite determinant winding
then counts physical eigenvalues with multiplicity.  No such numerical
enclosure is claimed by the construction alone.

\subsection{Endpoint expansions and physical tails}

\begin{proof}[Proof of Theorem~\ref{v12:thm:boundary-tail}]
For $P_\lambda u=0$, set $f=-\Delta_eu$.  The displayed kernel implies,
for every integer $j\ge0$,
\[
 |f^{(j)}(t)|\le C_j\|u\|_2\quad(0\le t\le2),\qquad
 |f^{(j)}(t)|\le C_j\|u\|_2t^{-1-j}\quad(t\ge2).
\]
On bounded rectangles, derivatives of divided differences are bounded
by cutoff derivatives; corner cancellation and separated denominators
handle the other entries.  At infinity one differentiates rational
tails.  These are estimates on kernels applied to $u\in L^2$, so no
regularity of $u$ has been assumed.  Consequently
$\widehat f=\mathcal Mf$ is holomorphic in
$|\operatorname{Im}\omega|<1/2$ and rapidly decreases horizontally on
every closed smaller strip: in logarithmic coordinates,
$e^{x/2}f(e^x)$ and all its derivatives decay exponentially at both
ends, allowing repeated integration by parts.

Since $M_\lambda Eu=f$, Mellin inversion gives
\[
 Eu(t)=\frac1{\sqrt{2\pi}}\int_{\mathbb R}
 t^{-1/2+i\omega}m_{\lambda,\theta}(\omega)^{-1}
                     \widehat f(\omega)\,d\omega.
\]
The determinant in \eqref{v12:eq:endpoint-symbol} has exactly two simple
zeros in this strip, at $\pm i\beta$.  Away from these poles the inverse has bounded horizontal limits.
Shifting down to $\operatorname{Im}\omega=-a$ crosses only $-i\beta$, giving
\[
 v=-i\sqrt{2\pi}\operatorname*{Res}_{\omega=-i\beta}
                  (m_{\lambda,\theta}^{-1}\widehat f).
\]
The remaining integral is $O(t^{-1/2+a})$ and may be differentiated
any fixed number of times.  Since $\chi=1$ near zero, this proves
\eqref{v12:eq:endpoint-expansion}.  A simple determinant zero has matrix
rank three, so the residue lies in the stated one-dimensional kernel;
it may vanish.

The interior error has smooth derivatives and rational tails.
From $N_\lambda\mathcal Iu=-\Delta_i u$ and
\eqref{v12:eq:interior-inverse}, the Fourier transform of $\mathcal Iu$
decays rapidly at high frequency and is locally $L^2$ at zero.
Thus $u$ is smooth away from its three endpoints.

A cutoff of $t^{\nu-1}$ at zero has Fourier asymptotic
$\Gamma(\nu)e^{-i\operatorname{sgn}(T)\pi\nu/2}
(2\pi|T|)^{-\nu}+O(|T|^{-1})$, by the regularized Gamma integral.
A remainder $r(t)=O(t^{\mu-1})$, $r'(t)=O(t^{\mu-2})$, $0<\mu<1$,
has Fourier transform $O(|T|^{-\mu})$: split at $|T|^{-1}$ and integrate
the upper part by parts.  Applying these formulas to the four endpoints
with $\mu=1/2+a$ gives \eqref{v12:eq:physical-tail}, with the phases
$e^{2\pi iT}$ and $e^{-2\pi iT}$ from endpoints $-1$ and $1$.

After squaring, the nonoscillating part integrates to
$L^{1-2\nu}/(2\nu-1)$; nonzero integer-frequency cross terms are
$O(L^{-2\nu})$ by integration by parts.  The remainder cross term is
$O(L^{1-\nu-\mu})=O(L^{-\beta-a})$, proving
\eqref{v12:eq:tail-mass}.  At $\omega=-i\beta$, the kernel equations
and $|q|=|r|=c=\sqrt{(1-\lambda^2)/2}$ give
\[
 |v_4|=|v_1|,\qquad
 |v_2|=|v_3|=
          \frac{\sqrt2|\lambda|}{\sqrt{1-\lambda^2}}|v_1|.
\]
Thus a nonzero residue makes both tail constants strictly positive,
even if the zero-frequency amplitude cancels.
\end{proof}

\begin{proof}[Proof of Corollary~\ref{v12:cor:gap-cluster-rate}]
Theorem~\ref{v13:thm:leading-interaction}, proved independently from the
boundary tails and the qualitative gap limit, gives the stronger
$\lambda_j=\alpha+L^{-\gamma(\alpha)}\mu_j+o(L^{-\gamma(\alpha)})$
for both boundaries and $\lambda=\alpha+o(L^{-\gamma(\alpha)})$
for one boundary. Since $L=m+h-l$, these imply the stated rate,
including each centered parity sector.
\end{proof}

